%%%%%%%%%%%%%%%%%%%%%%%%%%%%%%%%%%%%%%%%%%%%%%%%%%%%%%%%%%%%%%%%%%%%%%%%%%%%
%% Section 4: Simple Examples & Mnemonic Rules
%%%%%%%%%%%%%%%%%%%%%%%%%%%%%%%%%%%%%%%%%%%%%%%%%%%%%%%%%%%%%%%%%%%%%%%%%%%%

\begin{frame}{Exact vs Numerical Solutions}
\textbf{Analytical solution}

Closed-form expression.

Example:
\[
y' = y \implies y = Ce^x
\]

\textbf{Numerical methods}

When closed forms are impossible, compute solutions point-by-point:
\begin{itemize}
\item \textbf{Euler method} --- simplest: step forward using the slope
\item \textbf{Runge--Kutta} --- evaluate slope at multiple sub-steps for accuracy
\item \textbf{Finite difference methods} --- replace derivatives with discrete differences (PDEs)
\item \textbf{Finite element methods} --- approximate solutions on a mesh (complex geometries)
\end{itemize}

\vspace{0.3cm}
\textit{Most real-world differential equations have no closed-form solution. Numerical methods are not a fallback --- they are the primary tool.}
\end{frame}


\begin{frame}{Solving the IVP: Unambiguous Future}
\textbf{Problem:} $y' = y$, with the initial state $y(0) = 1$.

\textbf{Step 1: Find the General Solution} \\
This is a separable equation. We separate variables:
$$\frac{dy}{y} = dx$$

Integrate both sides:
$$\ln|y| = x + C \implies y(x) = Ae^x$$
\textit{(This represents a family of infinitely many curves.)}

\textbf{Step 2: Apply the Initial Condition} \\
Pin down the exact curve by plugging in our starting point $(0, 1)$:
$$1 = Ae^0 \implies A = 1$$

\textbf{Final Unique Solution:}
$y(x) = e^x$
\end{frame}

\begin{frame}{Solving the BVP: The Uniqueness Trap}
\textbf{Problem:} $y'' + y = 0,$ bounded at $x=0$ and $x=\pi.$

\textbf{Step 1: Find the General Solution} \\
The characteristic equation $r^2 + 1 = 0$ gives roots $r = \pm i$, yielding:
$$y(x) = c_1 \cos(x) + c_2 \sin(x)$$

\textbf{Step 2: Apply the First Boundary ($y(0) = 0$)}
$$0 = c_1 \cos(0) + c_2 \sin(0) \implies c_1 = 0$$
Our solution reduces to $y(x) = c_2 \sin(x)$.

\textbf{Step 3: The Second Boundary Dictates Solvability}
\begin{itemize}
    \item \textbf{If} $y(\pi) = 0$: $0 = c_2 \sin(\pi) \implies 0 = 0$. $c_2$ can be anything! \\ \textit{Result: Infinitely many solutions (e.g., any standing wave on a string).}
    \item \textbf{If} $y(\pi) = 1$: $1 = c_2 \sin(\pi) \implies 1 = 0$. Contradiction! \\ \textit{Result: Absolutely zero solutions.}
\end{itemize}
\end{frame}


\begin{frame}{Mnemonic Tricks \& Solution Recipes}
\textbf{1. Separation of variables} --- ``gather and integrate'' \\
If $y' = g(x)h(y)$, then $\displaystyle\int \frac{dy}{h(y)} = \int g(x)\,dx$.

\vspace{0.3cm}
\textbf{2. Integrating factor} --- for $y' + p(x)y = q(x)$ \\
Multiply both sides by $\mu(x) = e^{\int p(x)\,dx}$. The left side becomes $\frac{d}{dx}[\mu y]$.

\vspace{0.3cm}
\textbf{3. Characteristic equation} --- for constant-coefficient linear ODEs \\
Replace $y^{(k)}$ with $r^k$. Solve the polynomial: \\
\quad real roots $\to$ exponentials, \quad complex roots $\to$ sines/cosines, \quad repeated roots $\to$ multiply by $x$.

\vspace{0.3cm}
\textbf{4. The Ansatz philosophy} --- ``guess and check'' \\
If the driving term is $e^{ax}$, guess $y_p = Ae^{ax}$. \\
If the driving term is $\sin(\omega x)$, guess $y_p = A\sin(\omega x) + B\cos(\omega x)$. \\
Plug in, match coefficients, done.

\vspace{0.3cm}
\textit{Key insight: these are not separate tricks --- they are all consequences of linearity and the structure of the exponential function.}
\end{frame}
